\section{방법: 증명된 회전과 구성적 확장}
\label{sec:method}

정리~\ref{thm:optimal_hamiltonian}이 직접 정당화하는 것은 \emph{회전 선택}이다. 즉, Pre-RoPE 공분산의 PCA 기저를 사용해 pre-RoPE 키를 회전시키는 것이 Class~C 내 MSE 기준 최적이다. 반면 비트 배분(WF), residual tail 보존, Lloyd-Max 계열 양자화기는 이 증명 위에 올리는 \emph{구성적 확장}이며, 별도 실험 검증이 필요하다. 따라서 본 절은 (i) 증명된 PCA 기반 기본 파이프라인과 (ii) 실용 개선 후보를 분리해 기술한다.

\begin{algorithm}[t]
\caption{Pre-RoPE PCA 기반 기본 양자화 파이프라인}
\label{alg:pipeline}
\begin{algorithmic}[1]
\REQUIRE 캘리브레이션 키 $\{k^{(l,h)}_{\mathrm{pre}}\}$ (각 레이어 $l$, KV 헤드 $h$)
\REQUIRE 평균 비트 예산 $B$
\STATE \textbf{[오프라인: 캘리브레이션]}
\FOR{각 레이어 $l$, KV 헤드 $h$}
  \STATE $\Sigma_0^{(l,h)} \leftarrow \frac{1}{n}\sum_i (k_i - \bar{k})(k_i - \bar{k})^\top$ \COMMENT{중심화 공분산}
  \STATE $V_0^{(l,h)}, \{\lambda_j\} \leftarrow \mathrm{eigh}(\Sigma_0^{(l,h)})$ \COMMENT{고유분해}
\ENDFOR
\STATE \textbf{[온라인: 추론 시 k\_proj 후처리]}
\FOR{각 토큰의 pre-RoPE 키 $k_{\mathrm{pre}}$}
  \STATE $z \leftarrow (V_0^{(l,h)})^\top k_{\mathrm{pre}}$ \COMMENT{PCA 좌표 회전}
  \STATE $\hat{z}_j \leftarrow \mathrm{MinMaxQuant}(z_j, B)$ \quad $\forall j$ \COMMENT{기본선: 동일 비트 비대칭 양자화}
  \STATE $\hat{k}_{\mathrm{pre}} \leftarrow V_0^{(l,h)} \hat{z}$ \COMMENT{역회전}
  \STATE RoPE를 $\hat{k}_{\mathrm{pre}}$에 적용 \COMMENT{정상 어텐션 진행}
\ENDFOR
\end{algorithmic}
\end{algorithm}

\paragraph{캘리브레이션.} WikiText-2 train split에서 2048 토큰을 사용하여 각 KV 헤드별로 Pre-RoPE 키 벡터를 추출하고, 중심화 공분산의 고유분해를 수행한다. KVTC와 달리 \emph{헤드별} 독립 PCA를 계산한다 (정리~\ref{thm:optimal_hamiltonian}이 헤드별 최적화를 요구).

\paragraph{양자화.} 기본선은 채널별 비대칭 min-max 균일 양자화다: $\hat{z}_j = \mathrm{round}((z_j - z_{\min}) / s) \cdot s + z_{\min}$, 여기서 $s = (z_{\max} - z_{\min}) / (2^B - 1)$. 이것은 \texttt{k\_proj} 모듈의 forward hook으로 구현되어 모델 코드 수정 없이 적용된다.

\paragraph{구성적 확장 1: Query-Weighted Water-Filling (QW-WF).}
PCA 회전이 MSE-최적임은 정리로 증명되었다. 그 위의 비트 배분은 별도의 실용 설계이며, 본 논문은 다음의 query-weighted 배분을 제안한다:
\begin{equation}
b_j = B + \frac{1}{2}\log_2\frac{\lambda_j^{(K)} \cdot \sigma_j^{(Q)}}{\overline{\lambda^{(K)} \cdot \sigma^{(Q)}}}, \qquad b_j \geq 2,
\label{eq:qwwf}
\end{equation}
여기서 $\lambda_j^{(K)}$는 PCA 고유값, $\sigma_j^{(Q)} = \sqrt{(V_0^\top \Sigma_Q V_0)_{jj}}$는 PCA 좌표에서의 query 가중치다. 이 배분은 ``key 분산이 크고 동시에 query가 주목하는 차원''에 더 많은 비트를 할당한다.

\textbf{배경.} 명제~\ref{prop:attn_weighted}에서 attention-weighted 최적 회전은 $\Sigma_Q^{1/2} \Sigma_K \Sigma_Q^{1/2}$의 고유벡터를 요구하지만, 이것은 실험에서 PCA보다 악화되었다 ($\kappa(\Sigma_Q) \sim 10^4$로 인한 수치 불안정 + 분산 균일화). 반면, PCA 회전은 유지하면서 \emph{비트 배분에만} query 정보를 반영하는 QW-WF는 3모델에서 10--33\% PPL 개선을 보인다 (Section~\ref{sec:exp_qwwf}).

\textbf{PCA-Query 정렬.} 이 접근이 가능한 이유는, trained transformer에서 $\Sigma_K$와 $\Sigma_Q$의 상위 고유벡터가 자연적으로 정렬되어 있기 때문이다 (3모델 평균 주각 0.6--2.5°). 따라서 PCA 좌표에서 $\sigma_j^{(Q)}$를 계산하는 것이 의미 있다.

\paragraph{구성적 확장 2: residual tail 및 비균일 양자화}. 최근 토큰 tail을 FP16으로 남기거나 Lloyd-Max류 비균일 codebook을 추가하는 방법 역시 실용 개선 시도에 해당한다. 이들은 PPL이나 retrieval에서 유리할 수 있지만, 정리~\ref{thm:optimal_hamiltonian}의 증명 범위에는 포함되지 않는다.

\paragraph{계산 비용.} 캘리브레이션은 $O(n d_{\mathrm{head}}^2)$ (고유분해 1회), 추론 시 $O(d_{\mathrm{head}}^2)$ (행렬-벡터 곱 2회)이다. 무작위 회전 계열은 캘리브레이션이 불필요하지만, 추론 비용은 동일한 차수다.
